\section{Project Plan}

To ensure the systematic development and evaluation of the proposed dual-objective 
framework within the allotted research timeframe, the project will be executed using 
structured work packages, defined milestones, and proactive risk management protocols.

\subsection{Key Tasks}
The research methodology is partitioned into five sequential and logically dependent 
tasks.

\begin{enumerate}

\item \textbf{Infrastructure and Environment Provisioning}

\textbf{Objectives:} Establish a deterministic, hardware-agnostic execution environment.

\textbf{Tasks:}
\begin{itemize}
\item Deploy the base open-weights model (e.g., Llama-3-8B) on the target GPU 
infrastructure utilizing the vLLM serving engine.
\item Establish baseline latency metrics for unconstrained generation.
\end{itemize}

\item \textbf{Schema Compiler and Syntax Oracle Engineering}

\textbf{Objectives:} Develop the C++ backend for the symbolic constraint engine.

\textbf{Tasks:}
\begin{itemize}
\item Construct the Schema Compiler to parse the specification language for JSON Schemas.
\item Implement the FSM generator for the CFG mask.
\item Develop the topological-sort engine for evaluating the Cross-Field 
Dependency Graph ($G$).
\item Compile the \texttt{pybind11} bridges and verify inter-operatability.
\end{itemize}

\item \textbf{Cascaded Verification and Cache Rollback Mechanism}

\textbf{Objectives:} Integrate the Syntax Oracle into the autoregressive decoding loop.

\textbf{Tasks:}
\begin{itemize}
\item Subclass the \texttt{LogitsProcessor} to inject the token-level $O(1)$ FSM mask.
\item Implement the element-level trigger hooks.
\item Engineer the speculative cache rollback mechanism to 
truncate and regenerate upon Context-Sensitive Invariant (CSI) violations.
\end{itemize}

\item \textbf{Data Synthesis and Semantic Internalization}

\textbf{Objectives:} Align the neural policy to the target Domain-Specific Language (DSL).

\textbf{Tasks:}
\begin{itemize}
\item Utilize the Oracle to run coverage-guided random walks, 
synthesizing a mathematically perfect dataset $D$.
\item Execute Supervised Fine-Tuning (SFT) using Low-Rank 
Adaptation (LoRA) with optimized rank capacity ($r=16, \alpha=32$).
\end{itemize}

\item \textbf{Evaluation, Benchmarking, and Writing}

\textbf{Objectives:} Quantify the architectural synergy and draft the final dissertation.

\textbf{Tasks:}
\begin{itemize}
\item Measure the Entropy-Normalized Oracle Intervention Rate (OIR) 
across unaligned versus LoRA-aligned models.
\item Benchmark system throughput (tokens/second).
\item Finalize the manuscript/thesis.
\end{itemize}
\end{enumerate}

\subsection{Project Constraints}
To ensure the feasibility and scientific reproducibility of the framework, several boundary conditions are established:

\begin{itemize}
    \item  
    The framework will be evaluated against a restricted subset of a production DSL 
    (e.g., OpenAPI 3.0).

    \item  
    The training phase is computationally bounded by available GPU VRAM. 
    The use of PEFT (LoRA)~\cite{hu2021lora} specifically mitigates this 
    constraint, allowing adaptation to occur on a single consumer-grade 
    GPU or via standard university-provisioned SLURM clusters.

    \item  The Syntax Oracle is strictly coupled to the specific Byte-Pair Encoding 
    (BPE) vocabulary of the base model (e.g., Llama-3). 
    Any modification to the model's tokenizer requires a full recompilation of 
    the FSM transition table.
    
    \item  
    To maintain deterministic performance, the Semantic Evaluator is 
    constrained to \textit{local state invariants} present within the 
    current generation context. Verification against external, asynchronous 
    data sources (e.g., live database lookups) is outside the current research scope.

    \item  The speculative rollback mechanism is limited by the transformer's 
    effective context window. To prevent infinite loops 
    or "generation deadlocks," the system will impose a maximum retry threshold 
    for any single DSL element.
\end{itemize}